\part{Part IV: 基于机器学习的笔记热度分类预测模型}
\chapterimage{chapter_head_2.pdf} % Chapter heading image

\chapter{基准分类模型}

\section{预测问题描述}
在小红书内容传播效果研究中，一个核心任务是根据笔记的多维度特征（如内容、形式、发布者、时间、标签等）预测其是否具有高互动潜力。我们将此任务定义为二分类问题：给定一篇笔记的特征向量，模型需要预测其属于“高热度”（如互动量位于前20\%的爆款）还是“普通热度”类别。这比直接预测连续互动量更稳健，能有效处理极端值并服务于实际的二元决策（如是否加大推广）。通过构建和比较多种经典分类模型，我们可以识别出对预测贡献最大的特征，并理解各因素与热度之间的复杂关系，为内容优化提供可解释的洞见。

\section{逻辑回归}
逻辑回归（Logistic Regression）是解决二分类问题最基础的线性模型，以其简单、高效和高可解释性著称。尽管其假设特征与对数几率间存在线性关系，限制了其捕捉复杂非线性模式的能力，但作为基准模型，它能快速提供特征重要性的初步排序（通过系数大小和方向），帮助我们理解各因素对成为爆款的边际贡献。在处理小红书笔记分类时，逻辑回归可首先用于验证特征工程的合理性，并为后续复杂模型提供性能比较的基准。
\begin{itemize}
	\item \textbf{模型原理：} 通过Sigmoid函数将特征的线性组合映射到(0,1)区间，输出样本属于正类（高热度）的概率。决策边界为线性超平面。
	\item \textbf{优势：} 计算速度快，模型参数（系数）直接可解释，能直观看出各特征对结果的正/负向影响及相对强度。
	\item \textbf{局限：} 无法自动捕捉特征间复杂的交互作用和非线性关系，对数据线性可分性要求较高。
\end{itemize}

\section{支持向量机}
支持向量机（Support Vector Machine, SVM）通过寻找最大化分类间隔的最优超平面来实现分类，在处理高维、非线性问题上表现优异。对于小红书笔记数据，其高维特征（如文本向量、多标签）和潜在的非线性关系（如图片数量与热度的倒U型关系）使SVM成为一个强有力的基准模型。
\begin{itemize}
	\item \textbf{核技巧：} 通过引入高斯径向基函数（RBF）等核函数，SVM能将原始特征空间映射到更高维度，从而在变换后的空间中找到一个线性分界超平面，有效处理原始空间中的非线性模式。
	\item \textbf{优势：} 在高维空间中表现良好，泛化能力强，尤其适用于样本量不是特别巨大的情况。通过调整正则化参数C和核函数参数，可以灵活控制模型复杂度。
	\item \textbf{应用：} 在小红书笔记分类中，SVM可以学习到内容质量、发布者影响力、发布时间等特征之间复杂的相互作用，以区分爆款与普通笔记。
\end{itemize}

\section{集成学习模型}
集成学习通过结合多个基学习器的预测结果，获得比单一模型更稳定、更准确的强学习器，非常适合处理小红书笔记分类中特征多样、模式复杂的问题。

\subsection{随机森林}
随机森林（Random Forest）通过自助采样（Bootstrap）构建多棵决策树，并综合所有树的投票结果进行分类。其引入的随机特征选择有效降低了模型的方差，防止过拟合。
\begin{itemize}
	\item \textbf{优势：} 能处理高维数据，对缺失值不敏感，无需复杂的数据预处理。模型可提供特征重要性评分，有助于我们识别影响热度的关键因素（如封面图质量、特定话题标签）。
	\item \textbf{应用：} 在小红书场景下，随机森林可以同时评估文本、视觉、上下文等数百个特征，并给出重要性排序，例如可能发现“发布时段”和“标题情感”的交互作用对热度预测至关重要。
\end{itemize}

\subsection{梯度提升决策树}
梯度提升决策树（Gradient Boosting Decision Trees, GBDT）及其高效实现（如XGBoost, LightGBM）以串行方式迭代训练决策树，每一棵新树都致力于纠正前一棵树的残差。
\begin{itemize}
	\item \textbf{优势：} 预测精度通常高于随机森林，能自动处理特征的非线性关系和交互作用，且通过正则化控制模型复杂度。XGBoost和LightGBM在计算效率和内存使用上进行了大量优化，适合处理大规模数据集。
	\item \textbf{应用：} 对于拥有海量笔记数据的小红书研究，LightGBM能快速训练并高效处理类别型特征（如发布城市、星期几）。其输出的特征重要性（基于分裂增益或覆盖度）能深度揭示各因素对笔记热度的贡献，例如可能精确量化“使用特定热门标签”对互动量的提升幅度。
\end{itemize}

\section{深度学习模型}
深度学习模型能自动从原始或初级特征中学习高级表示，特别适合处理小红书笔记中的非结构化数据（如文本、图像）。

\subsection{多层感知机}
多层感知机（Multilayer Perceptron, MLP）是最基本的深度学习模型，由全连接层和非线性激活函数堆叠而成。它可以作为深度学习方法的入门基准，用于验证经过特征工程的结构化数据是否包含足够预测信息。
\begin{itemize}
	\item \textbf{优势：} 能够学习特征间的复杂非线性组合，相比逻辑回归具有更强的表示能力。
	\item \textbf{应用：} 可将经过预处理和特征工程的所有数值和类别特征（经过编码后）输入MLP，作为与其他更复杂深度学习模型对比的基线。
\end{itemize}

\subsection{卷积神经网络}
卷积神经网络（Convolutional Neural Network, CNN）在图像处理领域取得巨大成功，但其局部感知和参数共享的思想也可应用于一维序列数据。
\begin{itemize}
	\item \textbf{应用：} 对于文本数据，可将词嵌入序列视为一维“图像”，使用一维CNN来提取短语级别的局部语义特征。对于图像特征，可以直接使用预训练的CNN（如ResNet）提取的深度特征向量作为输入的一部分。CNN能有效捕捉如“标题中的特定关键词组合”或“封面图中同时出现人脸和产品”等局部模式对热度的信号。
\end{itemize}

\subsection{循环神经网络}
循环神经网络（Recurrent Neural Network, RNN）及其变体长短期记忆网络（LSTM）专门为序列数据设计，能捕捉时间或顺序上的依赖关系。
\begin{itemize}
	\item \textbf{应用：} 在小红书笔记分类中，LSTM可用于建模文本的序列信息（如正文的语义流），或建模用户历史行为序列（如发布者过去笔记的互动趋势）以捕捉其内容质量的时序演变。然而，对于单篇笔记的静态特征分类，LSTM的优势可能不如CNN或MLP明显，除非考虑时间上下文（如当前热点话题）。
\end{itemize}

\section{实证设置与基准结果}
为将上述模型由方案说明落实为可复现实验，本文读取工作区三个业务数据压缩包中的笔记详情表，合并后共得到323,598条原始记录。按笔记链接（无链接时按编号、标题和博主）去重并剔除互动指标缺失或为负的记录后，保留199,473条有效笔记。互动量定义为“点赞数+收藏数+$4\times$评论数”，达到全样本第80百分位（2,134）的笔记记为高热度，正类比例为20.01\%。建模输入仅使用发布前或内容本身可获得的10项特征，包括粉丝量对数、标题和正文长度、标签数、感叹号数、问号数、Emoji数、情感得分、星期和是否标注发布地；点赞、收藏和评论没有作为输入，以避免目标泄漏。

考虑计算成本，从有效数据中按类别分层抽取60,000条用于比较，再按8:2分层随机划分训练集与测试集。数值缺失以训练集的中位数填补；逻辑回归、线性SVM和MLP同时进行标准化。逻辑回归与SVM使用类别权重平衡，随机森林使用300棵树，梯度提升采用直方图梯度提升分类器，MLP包含64和32个神经元的两个隐藏层并使用早停。表\ref{tab:chapter10-baselines}报告固定测试集结果。鉴于正类仅占约20\%，模型选择以PR-AUC为主，ROC-AUC、F1、精确率和召回率为辅。

\begin{table}[htbp]
\centering
\caption{第10章基准分类模型测试集结果}
\label{tab:chapter10-baselines}
\begin{tabular}{lccccc}
\toprule
模型 & ROC-AUC & PR-AUC & F1 & 精确率 & 召回率 \\
\midrule
逻辑回归 & 0.700 & 0.405 & 0.414 & 0.303 & 0.649 \\
线性SVM & 0.700 & 0.406 & 0.413 & 0.302 & 0.651 \\
随机森林 & 0.725 & 0.424 & 0.373 & 0.479 & 0.305 \\
梯度提升树 & \textbf{0.736} & \textbf{0.440} & 0.263 & \textbf{0.600} & 0.169 \\
多层感知机 & 0.718 & 0.423 & 0.299 & 0.562 & 0.203 \\
\bottomrule
\end{tabular}
\end{table}

梯度提升树的排序能力最佳（PR-AUC=0.440，ROC-AUC=0.736），说明结构化特征与热度之间存在明显的非线性关系；随机森林和MLP次之。逻辑回归与线性SVM的PR-AUC接近，但在类别加权后召回率均约为0.65，更适合“尽量不漏掉潜在爆款”的初筛场景。梯度提升树默认阈值下精确率最高而召回率较低，若用于候选内容筛选，应在验证集上按业务成本降低分类阈值，而不能只依据默认0.5阈值下的F1作结论。

本轮完成了10.1至10.5.1的可运行实证。10.5.2的CNN与10.5.3的RNN/LSTM未给出数值结果：现有压缩包没有与每条笔记稳定对应的原始图片张量或发布者连续历史序列，直接把静态汇总特征重塑为“图像”或伪序列会制造不具业务含义的结构。虽然正文文本可形成词序列，但在没有独立时间外测试和文本模型调参预算的条件下，其结果不能与本节结构化基准作公平比较。因此本研究保留CNN与RNN的模型设计说明，待补齐图片、严格的用户历史序列和时间切分后再实施，避免报告不可复现或误导性的深度模型成绩。

\chapter{针对小红书数据特点的模型改进方向}

在基准模型的基础上，针对小红书笔记数据高维、稀疏、类别不平衡及多模态融合的特点，我们将探索以下改进方向以提升分类性能。

\section{处理类别不平衡}
小红书笔记热度分布高度不均衡，绝大多数笔记互动平平，仅少数成为爆款。直接使用原始数据训练会使模型严重偏向多数类。我们将采用多种策略：
\begin{itemize}
	\item \textbf{数据层重采样：} 使用SMOTE（Synthetic Minority Oversampling Technique）为少数类（高热度笔记）合成新样本，或使用欠采样技术（如NearMiss）减少多数类样本，或在模型训练中结合过采样与欠采样（如SMOTEENN）。
	\item \textbf{算法层代价敏感学习：} 在训练模型（如逻辑回归、SVM、决策树）时，为少数类样本设置更高的误分类代价（如调整`class\_weight`参数），迫使模型更关注爆款样本的学习。
	\item \textbf{集成方法：} 使用EasyEnsemble或BalanceCascade等集成方法，通过构建多个平衡的子训练集来训练基分类器，然后集成结果。
\end{itemize}

\section{多模态特征融合}
小红书笔记包含文本、图像、视频、上下文等多种模态信息。如何有效融合这些异构特征是提升模型性能的关键。
\begin{itemize}
	\item \textbf{早期融合：} 将不同模态的特征在输入层进行拼接，形成一个大特征向量，然后输入到传统机器学习模型（如XGBoost）或深度学习模型（如MLP）中。这种方法简单，但可能忽略模态间的深层交互。
	\item \textbf{晚期融合：} 为每个模态单独训练一个子模型（如用BERT处理文本，用ResNet处理图像），然后将各子模型的预测概率进行加权平均或作为元特征输入到一个融合模型（如逻辑回归）中进行最终决策。
	\item \textbf{中级融合与注意力机制：} 设计深度学习架构，分别用不同的子网络处理不同模态的数据，在中间层进行特征交互与融合。引入跨模态注意力机制（Cross-modal Attention），让模型动态地决定在做出分类决策时，哪些模态的哪些部分信息更重要。例如，对于一篇美妆教程笔记，模型可能更关注文本中的步骤描述和图片中的效果对比。
\end{itemize}

\section{引入图神经网络建模社交传播}
笔记的热度不仅取决于其自身质量，还与发布者的社交网络和社区互动有关。虽然数据受限，但我们可以尝试构建简化的图结构。
\begin{itemize}
	\item \textbf{构建笔记-用户二部图：} 节点分为笔记和用户两类。边表示用户与笔记的交互（点赞、收藏、评论）。通过图神经网络（如GCN、GraphSAGE）学习节点嵌入，笔记节点的嵌入可以融合其内容特征和其在图中的传播结构信息（如被哪些高影响力用户互动过），从而提升分类精度。
	\item \textbf{构建用户相似图：} 基于用户兴趣或行为相似性构建用户之间的边。通过学习用户节点的嵌入，可以丰富发布者特征，更好地表征其粉丝群体的潜在偏好。
\end{itemize}

\section{利用预训练语言与视觉模型}
迁移学习能大幅提升文本和图像特征的表示能力。
\begin{itemize}
	\item \textbf{文本特征：} 使用在大规模语料上预训练的中文BERT、RoBERTa等模型，对笔记标题和正文进行编码，获取包含丰富语义信息的上下文向量表示，替代传统的TF-IDF或Word2Vec特征。
	\item \textbf{图像特征：} 使用在ImageNet等大型数据集上预训练的CNN模型（如EfficientNet, CLIP）提取图像的通用视觉特征。CLIP模型因其对齐了图像和文本语义空间，特别适合理解小红书笔记中图文匹配的程度。
\end{itemize}

\chapter{模型评估与比较}

为了科学评估和比较不同分类模型的性能，我们将采用一系列评估指标，并重点关注模型在“高热度”少数类上的识别能力。

\section{评估指标}
鉴于数据集的不平衡性，准确率（Accuracy）不是可靠的评估指标。我们将主要使用以下指标：
\begin{itemize}
	\item \textbf{精确率 (Precision)：} 模型预测为爆款的笔记中，真正是爆款的比例。高精确率意味着推送或推广的误判成本低。
	\item \textbf{召回率 (Recall)：} 所有真实爆款笔记中，被模型正确找出的比例。高召回率意味着不错过潜在的爆款机会。
	\item \textbf{F1-Score：} 精确率和召回率的调和平均数，是综合衡量模型在正类上性能的单一指标。我们主要参考宏平均F1（Macro-F1），它平等看待每一类，对少数类性能变化更敏感。
	\item \textbf{ROC曲线与AUC：} 绘制真正例率（TPR/Recall）随假正例率（FPR）变化的曲线。曲线下面积（AUC）衡量模型整体区分类别的能力，AUC越接近1，模型性能越好。ROC-AUC对类别不平衡相对不敏感。
	\item \textbf{PR曲线与AUC：} 绘制精确率随召回率变化的曲线。在不平衡数据中，PR曲线比ROC曲线更能反映模型在正类上的真实性能。PR曲线下面积（PR-AUC）越高，表明模型在保证高精确率的同时也能获得高召回率的能力越强。
	\item \textbf{混淆矩阵：} 直观展示模型在四个类别（TP, FP, FN, TN）上的具体表现，便于分析错误类型。
\end{itemize}

\section{模型比较与选择流程}
我们将遵循以下严谨流程进行模型比较与选择：
\begin{enumerate}
	\item \textbf{数据划分：} 将数据集按时间顺序划分为训练集、验证集和测试集（如7:2:1），以模拟现实中的时间外预测，避免未来信息泄露。
	\item \textbf{交叉验证：} 在训练集上使用时序交叉验证（TimeSeriesSplit）进行超参数调优和模型选择，确保评估的稳健性。
	\item \textbf{基准线：} 设定一个简单的基准线（如预测所有笔记为普通类，或基于粉丝数的简单规则模型），所有复杂模型性能需显著优于该基准线。
	\item \textbf{综合评估：} 在验证集上，综合比较各模型的宏平均F1、ROC-AUC、PR-AUC等指标。优先选择在PR-AUC和召回率上表现突出的模型，因为本研究更关注尽可能多地识别出潜在爆款（高召回），同时控制误判率（高精确率）。
	\item \textbf{测试集验证：} 最终选定1-3个最佳模型，在独立的测试集上进行最终性能报告，确保结果的可靠性。
	\item \textbf{可解释性分析：} 对最优模型（如XGBoost、带注意力机制的神经网络）进行可解释性分析，使用SHAP值、特征重要性图、LIME等方法，量化并可视化各特征对预测结果的贡献，从而得出可操作的、用于提升笔记热度的策略建议。
\end{enumerate}
